% ============================================================
%  Práctica de Laboratorio 2 — Equipo ISA (Motor CC)
%  Control Automático — 2º Grado en ITT
%  Dpto. Ingeniería de Sistemas y Automática — Universidad de Sevilla
% ============================================================
\documentclass[12pt,a4paper]{article}
\usepackage[utf8]{inputenc}
\usepackage[T1]{fontenc}
\usepackage[spanish,es-nodecimaldot]{babel}
\usepackage{amsmath,amssymb}
\usepackage{graphicx}
\usepackage{tikz}
\usepackage{pgfplots}
\usepackage{booktabs}
\usepackage{array}
\usepackage{geometry}
\usepackage[colorlinks=true,linkcolor=blue!60!black,urlcolor=blue!60!black]{hyperref}
\usepackage{xcolor}
\usepackage{tcolorbox}
\usepackage{enumitem}
\usepackage{float}
\usepackage{caption}
\usepackage{subcaption}
\usepackage{fancyhdr}
\usepackage{multirow}
\usepackage{siunitx}

\usetikzlibrary{arrows.meta, positioning, fit, calc, shapes.geometric,
                shapes.misc, decorations.pathreplacing, backgrounds}
\pgfplotsset{compat=1.18}

\geometry{margin=2.5cm, top=3cm, bottom=3cm, headheight=14pt}

\definecolor{isablue}{RGB}{0,70,130}
\definecolor{isared}{RGB}{180,30,30}
\definecolor{lightgray}{RGB}{240,240,245}

\pagestyle{fancy}
\fancyhf{}
\lhead{\small\textcolor{isablue}{\textbf{Control Automático — ITT}}}
\rhead{\small\textcolor{isablue}{\textbf{Práctica 2 — Equipo ISA}}}
\cfoot{\small\thepage}
\renewcommand{\headrulewidth}{0.4pt}

\tcbuselibrary{skins,breakable}
\newtcolorbox{pregunta}[2][]{
  enhanced,
  colback=isared!8,
  colframe=isared,
  fonttitle=\bfseries\small,
  title={#2},
  #1
}
\newtcolorbox{notabox}{
  enhanced,
  colback=isablue!8,
  colframe=isablue,
  fonttitle=\bfseries\small,
  title={Nota},
}

% ---- Estilo bloques TikZ ----
\tikzset{
  sumblock/.style={draw, circle, fill=white, minimum size=0.6cm,
                   inner sep=0pt, font=\small},
  ctrlblock/.style={draw, fill=isablue!15, minimum width=1.8cm,
                    minimum height=0.8cm, align=center, font=\small,
                    rounded corners=2pt},
  plantblock/.style={draw, fill=isared!10, minimum width=2.4cm,
                     minimum height=0.8cm, align=center, font=\small,
                     rounded corners=2pt},
  arr/.style={-{Stealth[length=5pt]}, semithick}
}

% ===========================================================
\begin{document}
% ===========================================================

\begin{center}
  {\large Control Automático}\\[2pt]
  {\large 2\textsuperscript{o} Curso Grado en Ingeniería de las Tecnologías de la Telecomunicación}\\[10pt]
  \rule{\linewidth}{1.5pt}\\[6pt]
  {\LARGE\bfseries\color{isablue} Práctica de Laboratorio 2}\\[4pt]
  {\Large Control PID de un motor de corriente continua}\\[4pt]
  \rule{\linewidth}{1.5pt}\\[6pt]
  {\small Dpto.\ de Ingeniería de Sistemas y Automática — ISA\\
   Universidad de Sevilla \quad | \quad \url{ialvarado@us.es}}
\end{center}

\vspace{6pt}

% ============================================================
\section{Objetivos}
% ============================================================

Esta práctica continúa el trabajo iniciado en la práctica~1 con el equipo ISA.
Los objetivos son:

\begin{itemize}[noitemsep]
  \item Repetir la identificación del sistema y analizar la zona muerta del motor.
  \item Diseñar e implementar un controlador de \textbf{velocidad}:
        primero proporcional (P) y después proporcional-integral (PI).
  \item Diseñar e implementar un controlador de \textbf{posición}:
        proporcional-derivativo (PD) y proporcional-integral-derivativo (PID).
  \item Comprender la influencia de cada parámetro del controlador PID y el
        efecto de la zona muerta en el error en régimen permanente.
\end{itemize}

\begin{notabox}
  Para toda esta práctica puede suponer que el punto de funcionamiento en tensión
  es nulo ($V_0=0$), por lo que no es necesario el paso a variables incrementales.
\end{notabox}

% ============================================================
\section{Análisis del sistema}
% ============================================================

\subsection{Identificación (revisión)}

\begin{enumerate}
  \item Abra el modelo \texttt{MotorCC\_V2.slx} y configure el puerto serie
        haciendo doble clic en el bloque Motor\_CC.
  \item Partiendo del modelo de Simulink de la práctica~1 (bloque Motor\_CC con
        la salida de velocidad en rad/s, opcionalmente filtrada), aplique un
        escalón desde $u=0$ hasta $u=42$ ($V_{eq}\approx\SI{2.0}{V}$) y
        verifique que el modelo funciona correctamente.
\end{enumerate}

\begin{pregunta}{Pregunta \#1 — Función de transferencia}
  Repita la identificación del apartado~3 de la práctica~1 con el escalón
  $u=42$ ($V_{eq}\approx\SI{2.0}{V}$). Anote la función de transferencia
  $G(s)=K/(\tau s+1)$ con los valores de $K$ y $\tau$ obtenidos.
  Esta función de transferencia se utilizará en los apartados siguientes.
\end{pregunta}

\subsection{Zona muerta}

La zona muerta es un fenómeno no lineal: por debajo de cierta tensión de entrada,
el motor no se mueve porque la fricción estática del conjunto motor+reductora supera
el par motor. Los límites de zona muerta son distintos para cada sentido de giro.

\begin{enumerate}
  \item Configure el bloque \emph{Constant} con un valor inicial $u=5$.
        Ejecute la simulación y compruebe si el motor se mueve observando la señal
        de velocidad $\omega$ en el Scope.
  \item Vaya aumentando $u$ de en 1 en 1 hasta detectar el valor mínimo para el
        cual el motor gira. Anótelo como $u_{ZM}^{+}$.
  \item Repita con valores negativos para encontrar $u_{ZM}^{-}$.
\end{enumerate}

\begin{pregunta}{Pregunta \#2 — Zona muerta}
  \begin{enumerate}[label=\alph*)]
    \item ¿Cuáles son los límites de la zona muerta
          $u_{ZM}^{+}$ y $u_{ZM}^{-}$?
    \item Conviértalos a tensión equivalente: $V_{ZM}^{\pm}=u_{ZM}^{\pm}\cdot12/255$.
    \item ¿Es la zona muerta simétrica para ambos sentidos? ¿Por qué podría no serlo?
    \item ¿Qué implicación tiene la zona muerta para los controladores que diseñará
          a continuación?
  \end{enumerate}
\end{pregunta}

% ============================================================
\section{Control de la velocidad del motor}
% ============================================================

La figura~\ref{fig:vel_loop} muestra el esquema de control de velocidad en bucle
cerrado. La referencia $\omega_r$ se compara con la velocidad medida $\omega$;
el error $e=\omega_r-\omega$ es procesado por el controlador $C(s)$ para generar
la señal de control $u$.

\begin{figure}[H]
\centering
\begin{tikzpicture}[node distance=1.3cm]
  \node[sumblock] (sum) {\footnotesize$+$};
  \node[ctrlblock, right=1.4cm of sum] (ctrl) {$C(s)$};
  \node[plantblock, right=1.4cm of ctrl] (plant)
    {$G(s)=\dfrac{K}{\tau s+1}$};
  \node[right=1.4cm of plant] (out) {};

  % Referencias y salidas
  \node[left=1.4cm of sum] (ref) {$\omega_r$};
  \draw[arr] (ref) -- (sum);
  \draw[arr] (sum) -- node[above,font=\scriptsize]{$e$} (ctrl);
  \draw[arr] (ctrl) -- node[above,font=\scriptsize]{$u$} (plant);
  \draw[arr] (plant) -- node[above,font=\scriptsize]{$\omega$} (out.center);

  % Realimentación
  \draw[arr] (out.center) +(0.4,0) -- +(0.4,-1.0) -| node[near end,right,font=\scriptsize]{$-$} (sum.south);

  % Etiqueta
  \node[below=0.1cm of sum, font=\tiny, xshift=0.1cm] {$-$};
\end{tikzpicture}
\caption{Esquema de control de velocidad en bucle cerrado.}
\label{fig:vel_loop}
\end{figure}

\subsection{Controlador proporcional (P)}

\begin{enumerate}
  \item Modifique el modelo de Simulink para implementar un \textbf{controlador P}:
        añada un bloque \emph{Sum} que reste la velocidad medida a la referencia,
        conecte la salida al bloque Motor\_CC a través de un \emph{Gain}~$K_p$.
  \item Como referencia, use un \emph{Pulse Generator} con amplitud $\pm 20\;\si{rad/s}$
        y frecuencia $\SI{0.4}{Hz}$ (período $\SI{2.5}{s}$).
  \item Pruebe distintos valores de $K_p$ y analice el comportamiento del sistema
        (error en régimen permanente, sobreimpulso, oscilaciones).
  \item Introduzca perturbaciones presionando suavemente el eje del motor con
        un dedo y observe la reacción del controlador.
\end{enumerate}

\begin{pregunta}{Pregunta \#3 — Controlador proporcional}
  \begin{enumerate}[label=\alph*)]
    \item ¿Existe error en régimen permanente con el controlador~P?
          ¿A qué se debe? (Pista: recuerde el tipo del sistema en bucle cerrado.)
    \item ¿Cómo afecta aumentar $K_p$ al error y a la estabilidad?
    \item ¿Responde el sistema adecuadamente a las perturbaciones?
          ¿Qué valor de $K_p$ parece un compromiso razonable?
  \end{enumerate}
\end{pregunta}

\subsection{Diseño analítico del controlador PI}

El objetivo de diseño es conseguir que el sistema en bucle cerrado se comporte como
un sistema de primer orden con constante de tiempo $\tau_{cl}=\SI{0.1}{s}$ y ganancia
estática unitaria.

El controlador PI tiene la forma:
\begin{equation}
  C(s) = K_p\,\frac{T_i s + 1}{T_i s}
  \label{eq:PI}
\end{equation}

La función de transferencia en bucle cerrado con el controlador~(\ref{eq:PI}) y la
planta $G(s)=K/(\tau s+1)$ es:
\begin{equation}
  T(s) = \frac{C(s)\,G(s)}{1+C(s)\,G(s)}
\end{equation}

\textbf{Procedimiento de diseño (cancelación polo-cero):}

\begin{enumerate}
  \item \textbf{Cancelación:} Elija $T_i = \tau$ para cancelar el polo de la planta.
        Así $C(s)\,G(s) = K_p K / (T_i s)$.
  \item \textbf{Especificación:} La función de bucle cerrado resultante es:
        \begin{equation}
          T(s) = \frac{1}{\frac{T_i}{K_p K}\,s + 1}
          = \frac{1}{\tau_{cl}\,s+1}
          \quad\Rightarrow\quad
          \tau_{cl} = \frac{T_i}{K_p K}
          \label{eq:tau_cl}
        \end{equation}
  \item \textbf{Cálculo:} De (\ref{eq:tau_cl}) se obtiene
        $K_p = T_i / (K\,\tau_{cl}) = \tau / (K \cdot \SI{0.1}{s})$.
\end{enumerate}

\begin{enumerate}
  \item Calcule $K_p$ y $T_i$ con la función de transferencia identificada en \#1.
  \item Implemente el controlador PI en Simulink usando un bloque \emph{Transfer Fcn}
        con numerador $[K_p T_i,\;K_p]$ y denominador $[T_i,\;0]$,
        o bien con un bloque \emph{PID Controller} con $K_p$, $K_i=K_p/T_i$, $K_d=0$.
  \item Añada una saturación a la salida del controlador para limitar $u\in[-255,255]$.
  \item Ejecute la simulación con la referencia en escalones $\pm 20\;\si{rad/s}$
        a $\SI{0.4}{Hz}$.
\end{enumerate}

\begin{pregunta}{Pregunta \#4 — Control PI de velocidad}
  \begin{enumerate}[label=\alph*)]
    \item Anote los valores de $K_p$ y $T_i$ calculados.
    \item ¿Se elimina el error en régimen permanente? ¿Por qué el PI consigue esto?
    \item Si es necesario, ajuste $K_p$ manualmente para mejorar la respuesta
          (rapidez, sobreimpulso). Anote el valor final utilizado.
    \item Introduzca perturbaciones con el dedo. ¿Cómo responde el PI
          frente al controlador~P? Razone en base a la acción integral.
  \end{enumerate}
\end{pregunta}

% ============================================================
\section{Control de la posición del motor}
% ============================================================

Ahora la variable a controlar es la \textbf{posición angular} $\theta$ en radianes.
El esquema de control es idéntico al de la figura~\ref{fig:vel_loop} pero con
$\theta_r$ como referencia y $\theta$ como salida (figura~\ref{fig:pos_loop}).

\begin{figure}[H]
\centering
\begin{tikzpicture}[node distance=1.3cm]
  \node[sumblock] (sum) {\footnotesize$+$};
  \node[ctrlblock, right=1.4cm of sum] (ctrl) {$C(s)$};
  \node[plantblock, right=1.5cm of ctrl] (plant)
    {$G_\theta(s)=\dfrac{K}{\tau s^2+s}$};
  \node[right=1.4cm of plant] (out) {};

  \node[left=1.4cm of sum] (ref) {$\theta_r$};
  \draw[arr] (ref) -- (sum);
  \draw[arr] (sum) -- node[above,font=\scriptsize]{$e$} (ctrl);
  \draw[arr] (ctrl) -- node[above,font=\scriptsize]{$u$} (plant);
  \draw[arr] (plant) -- node[above,font=\scriptsize]{$\theta$} (out.center);
  \draw[arr] (out.center) +(0.4,0) -- +(0.4,-1.0) -| node[near end,right,font=\scriptsize]{$-$} (sum.south);
  \node[below=0.1cm of sum, font=\tiny, xshift=0.1cm] {$-$};
\end{tikzpicture}
\caption{Esquema de control de posición en bucle cerrado.
La planta $G_\theta(s)$ incluye el integrador implícito entre velocidad y posición.}
\label{fig:pos_loop}
\end{figure}

La función de transferencia de la posición respecto a la tensión es:
\begin{equation}
  G_\theta(s) = \frac{\Theta(s)}{V_{eq}(s)} = \frac{G(s)}{s} = \frac{K}{s(\tau s+1)}
\end{equation}
pues la posición es la integral de la velocidad.

\subsection{Diseño analítico del controlador PD}

El objetivo de diseño para posición es un sistema en bucle cerrado de primer orden
con $\tau_{cl}=\SI{0.02}{s}$ y ganancia estática unitaria.

El controlador PD (con la derivada filtrada para que sea implementable):
\begin{equation}
  C(s) = K_p\left(1 + T_d\,s\right)
  \label{eq:PD}
\end{equation}

\textbf{Procedimiento (cancelación polo-cero):}
\begin{enumerate}
  \item Elija $T_d = \tau$ para cancelar el polo de la planta.
        La función de bucle resulta $C(s)\,G_\theta(s) = K_p K / s$.
  \item La función de bucle cerrado es:
        \begin{equation}
          T(s) = \frac{K_p K / s}{1 + K_p K / s} = \frac{1}{\frac{1}{K_p K}s+1}
          = \frac{1}{\tau_{cl}\,s+1}
          \quad\Rightarrow\quad
          K_p = \frac{1}{K\,\tau_{cl}}
        \end{equation}
  \item Calcule $K_p$ con $\tau_{cl}=\SI{0.02}{s}$.
\end{enumerate}

\begin{notabox}
  La acción derivativa pura es no implementable. Use el bloque \emph{PID Controller}
  de Simulink con $K_d$ y un filtro de derivada con coeficiente $N$
  (el PD queda: $C(s) = K_p + K_p T_d s\,N/(s+N)$). Comience con $N=100$.
\end{notabox}

\begin{enumerate}
  \item Modifique el modelo de Simulink para controlar la posición:
        \begin{itemize}[noitemsep]
          \item Conecte la salida de posición $\theta$ [°] del bloque Motor\_CC,
                conviértala a rad (multiplicar por $\pi/180$) y réaliméntela.
          \item Como referencia, use un \emph{Pulse Generator} con amplitud
                $\pm 0.5\;\si{rad}$ y frecuencia $\SI{0.1}{Hz}$ (período $\SI{10}{s}$).
        \end{itemize}
  \item Implemente el controlador PD con $K_p$ calculado y $T_d=\tau$.
  \item Ejecute la simulación y analice el resultado.
\end{enumerate}

\begin{pregunta}{Pregunta \#5 — Control PD de posición}
  \begin{enumerate}[label=\alph*)]
    \item Anote los valores de $K_p$ y $T_d$ calculados.
    \item ¿Se obtiene el comportamiento de primer orden esperado con $\tau_{cl}=\SI{0.02}{s}$?
    \item ¿Existe error en régimen permanente? ¿A qué se debe?
          (Pista: piense en el efecto de la zona muerta.)
    \item Si es necesario, ajuste $K_p$ manualmente para mejorar la respuesta.
          ¿Qué compromiso aparece al aumentar $K_p$?
  \end{enumerate}
\end{pregunta}

\subsection{Control PID: reducción del error en régimen permanente}

El error en régimen permanente con control PD se debe principalmente a la zona
muerta: el controlador no puede generar la pequeña acción de control necesaria
para compensar la fricción estática. La acción integral del PID puede reducir
(aunque no eliminar completamente) este error.

\begin{enumerate}
  \item Parta del controlador PD anterior y añada acción integral:
        \begin{equation}
          C(s) = K_p\left(1 + \frac{1}{T_i s} + T_d s\right)
        \end{equation}
        Comience con un valor de $T_i$ grande (p.~ej., $T_i = \SI{1}{s}$) y
        vaya reduciéndolo mientras el sistema permanezca estable.
  \item Observe el efecto de la acción integral sobre el error en régimen
        permanente y sobre el sobreimpulso.
  \item Ajuste los tres parámetros $K_p$, $T_i$ y $T_d$ por prueba y error para
        obtener la mejor respuesta posible (mínimo error y mínimo sobreimpulso).
\end{enumerate}

\begin{pregunta}{Pregunta \#6 — Control PID de posición}
  \begin{enumerate}[label=\alph*)]
    \item Anote los parámetros finales del PID ($K_p$, $T_i$, $T_d$).
    \item ¿Se reduce el error en régimen permanente con la acción integral?
          ¿Lo elimina completamente? ¿Por qué?
    \item ¿Qué problema aparece si $T_i$ es demasiado pequeño?
    \item Razone si es posible diseñar un PID analíticamente cuando el sistema
          presenta zona muerta. ¿Qué herramientas de diseño serían necesarias?
    \item Compare las respuestas PD y PID en la misma gráfica.
          ¿Cuál preferiría para esta aplicación? Justifique la respuesta.
  \end{enumerate}
\end{pregunta}

% ============================================================
\section*{Hoja de respuestas}
% ============================================================
\addcontentsline{toc}{section}{Hoja de respuestas}

\begin{center}
{\large\bfseries Hoja de respuestas — Práctica 2}
\end{center}

Nombre: \underline{\hspace{8cm}}

\vspace{6pt}

\begin{tabular}{|p{5.5cm}|p{8.5cm}|}
\hline
\multicolumn{1}{|c|}{\textbf{Dato solicitado}} &
\multicolumn{1}{c|}{\textbf{Respuesta}} \\
\hline
\#1 Función de transferencia $G(s)$ &
\vspace{2.5cm} \\
\hline
\#2 Límites de zona muerta ($u_{ZM}^+$, $u_{ZM}^-$) &
\vspace{2.5cm} \\
\hline
\#3 Comportamiento controlador P ($K_p$ elegido) &
\vspace{2.5cm} \\
\hline
\#4 Control PI de velocidad ($K_p$, $T_i$) &
\vspace{2.5cm} \\
\hline
\#5 Control PD de posición ($K_p$, $T_d$) &
\vspace{2.5cm} \\
\hline
\#6 Control PID de posición ($K_p$, $T_i$, $T_d$) &
\vspace{3.0cm} \\
\hline
\end{tabular}

% ============================================================
\appendix
\section{Implementación del controlador PID en Simulink}
% ============================================================

Simulink dispone del bloque \emph{PID Controller} (en la librería
\texttt{Simulink/Continuous}) que implementa directamente el PID en la forma:
\begin{equation}
  u(t) = K_p\,e(t) + \frac{K_p}{T_i}\int_0^t e(\sigma)\,d\sigma
         + K_p\,T_d\,\frac{N\,e(t)}{1 + N\int_0^t e(\sigma)\,d\sigma\cdot\ldots}
\end{equation}
o de forma equivalente en frecuencia:
\begin{equation}
  C(s) = K_p + \frac{K_p}{T_i\,s} + K_p\,T_d\,\frac{N\,s}{s+N}
\end{equation}
donde $N$ es el coeficiente de filtrado de la derivada. Un valor típico de
$N\in[20,\,200]$ asegura que la derivada sea implementable y no amplifique
excesivamente el ruido de alta frecuencia.

Configure el bloque PID en MATLAB/Simulink de la siguiente manera:
\begin{enumerate}[noitemsep]
  \item Arrastre un bloque \emph{PID Controller} al diagrama.
  \item Haga doble clic y seleccione el tipo (P, PI, PD o PID).
  \item Introduzca los parámetros $K_p$ (proporcional), $T_i$ (integral),
        $T_d$ (derivativo) y $N$ (filtro de derivada).
  \item Añada siempre un bloque \emph{Saturation} a la salida del controlador
        limitando $u\in[-255,\,255]$ para no saturar el driver del motor.
\end{enumerate}

\subsection*{Anti-windup}

Cuando la señal de control satura, el integrador puede seguir acumulando error
(\emph{windup}), provocando sobretensiones al salir de la saturación.
Habilite el anti-windup en el bloque PID seleccionando el método
\emph{Back-calculation} con constante de rastreo $K_b = 1/T_i$.

\end{document}
